\textbf{COLUNI 2025 -- Questão 25}\par\medskip Considere as informações sobre urânio: os minerais que contêm urânio naturalmente apresentam principalmente os isótopos urânio-238 (U-238) e urânio-235 (U-235). O U-235 é o isótopo que pode sustentar uma reação nuclear em cadeia, mas é menos abundante, de acordo com a tabela abaixo.

\begin{center}
\begin{tabular}{|c|c|c|}
\hline
Símbolo do Isótopo & \% no mineral natural & Massa atômica \\
\hline
U-235 & 0,726 & 235 \\
U-238 & 99,274 & 238 \\
\hline
\end{tabular}
\end{center}

O processo de enriquecimento de urânio aumenta a concentração de U-235 na amostra. Um dos métodos mais comuns envolve a conversão do urânio em um gás, que é então colocado em um dispositivo que gira a altas velocidades. Durante este processo, o isótopo mais leve tende a se concentrar no centro, enquanto o mais pesado se acumula na periferia. Este procedimento é repetido várias vezes até atingir a concentração desejada de U-235. O resultado final é uma amostra com maior concentração de urânio (U-235) que a amostra inicial, conhecido como urânio enriquecido, com maior proporção de U-235, que é usada em reatores para produção de energia elétrica e armas nucleares.

Com base nas informações fornecidas e nos conhecimentos sobre ciências, assinale a afirmativa CORRETA:
\begin{enumerate}
\item O isótopo mais abundante no mineral possui 143 nêutrons em sua composição nuclear.
\item O método de separação para obtenção do urânio enriquecido é conhecido como centrifugação.
\item O urânio é um elemento químico do sétimo período da tabela periódica, pertencente aos ametais.
\item Os átomos de urânio-238 (U-238) e urânio-235 (U-235) são reconhecidos como isóbaros.
\end{enumerate}
